\documentclass[11pt]{article}
\usepackage[margin=1in]{geometry}
\usepackage{amsmath,amssymb,amsthm}
\usepackage[T1]{fontenc}
\usepackage{lmodern}
\usepackage{xcolor}
\usepackage{listings}
\usepackage{booktabs}
\usepackage{longtable}
\usepackage{array}
\usepackage[hidelinks]{hyperref}
\usepackage{url}

\newtheorem{theorem}{Theorem}
\newtheorem{lemma}[theorem]{Lemma}
\newtheorem{corollary}[theorem]{Corollary}
\newtheorem{proposition}[theorem]{Proposition}
\theoremstyle{definition}
\newtheorem{definition}[theorem]{Definition}
\theoremstyle{remark}
\newtheorem{remark}[theorem]{Remark}

% Code-like text (not checked by the name gate).
\newcommand{\lean}[1]{\texttt{#1}}
% A declaration of this development. Every \decl{...} is checked by namegate.sh
% (a declaration of that name, in the stated namespace, in Baranyai/*.lean).
\DeclareUrlCommand\decl{\urlstyle{tt}\def\UrlBreaks{\do\_\do\.}}
% A Mathlib declaration. Every \mdecl{...} is checked by namegate.sh against
% checks/mathlib_names.txt (names #check-ed against the pinned Mathlib).
\DeclareUrlCommand\mdecl{\urlstyle{tt}\def\UrlBreaks{\do\_\do\.}}
% version-2 commits: baranyai=1ab21623da9572cb10a8630a086c010b9e554269
\newcommand{\repo}{https://github.com/jpbald93/baranyai}
\newcommand{\commit}{051ae8c4063be0bb70f1726e9f60d0319dd37d91}
\newcommand{\shortcommit}{051ae8c}
% \src{File}{n}: link to line n of Baranyai/File.lean at \commit.
\newcommand{\src}[2]{\href{\repo/blob/\commit/Baranyai/#1.lean\#L#2}{\texttt{#1:#2}}}
% \file{path}: link to a file of the repository at \commit.
\newcommand{\file}[1]{\href{\repo/blob/\commit/#1}{\path{#1}}}
\newcommand{\NN}{\mathbb{N}}

\lstdefinelanguage{lean}{
  morekeywords={theorem,lemma,def,structure,where,by,fun,namespace,end,import,open},
  sensitive=true, morecomment=[l]{--}, morecomment=[s]{/-}{-/}}
\lstset{language=lean, basicstyle=\ttfamily\small, keywordstyle=\bfseries,
  columns=fullflexible, keepspaces=true, breaklines=true, frame=single, framerule=0.3pt,
  literate={ℕ}{{$\mathbb{N}$}}1 {≠}{{$\neq$}}1 {↔}{{$\leftrightarrow$}}2
    {→}{{$\to$}}1 {∀}{{$\forall$}}1 {∃}{{$\exists$}}1 {∧}{{$\wedge$}}1 {∈}{{$\in$}}1
    {≤}{{$\le$}}1 {∣}{{$\mid$}}1 {≃}{{$\simeq$}}1 {×}{{$\times$}}1 {α}{{$\alpha$}}1
    {β}{{$\beta$}}1 {γ}{{$\gamma$}}1 {⟨}{{$\langle$}}1 {⟩}{{$\rangle$}}1}

\title{Baranyai's theorem on factorisations of complete uniform hypergraphs,\\ formalised in
Lean~4 with Mathlib}
\author{Josh Bald\\ \small Independent researcher, Ontario, Canada\\ \small\texttt{jpbald93@gmail.com}}
\date{2 October 2026; version 2, 22 October 2026 (adds the Layer-14 and Layer-15 checks of
Section~\ref{sec:l1415})}

\emergencystretch=3em
\begin{document}
\maketitle

\begin{abstract}
For positive integers $n,k$, Baranyai's theorem (1975) says that if $k$ divides $n$, the $\binom nk$ $k$-element subsets of an
$n$-element set can be split into $\binom{n-1}{k-1}$ classes, each of which is a partition of the
$n$-set into $n/k$ blocks. For $k=2$ this is the classical 1-factorisation of the complete graph
$K_n$, $n$ even. We describe a Lean~4 formalisation, built on Mathlib, of the theorem for vertex
set $\{0,\dots,n-1\}$ (the formal colouring statement also admits $n=0$, where it holds only
vacuously), together with its transport to an arbitrary finite type (and to the
elements of a \lean{Finset}) and, for $n>0$, a partition form that spells out the disjointness,
covering and class-size properties of the colour classes: each class has $n/k$ blocks (a division
lemma that needs $k>0$), and the number of classes satisfies $\binom{n-1}{k-1}=\binom nk/(n/k)$ (an
identity that needs $n>0$, $k>0$ and $k\mid n$, and fails for $n=0$, $k=1$). The proof is Baranyai's vertex-by-vertex induction in the form
given by Brouwer and Schrijver; the integral rounding step is obtained from Mathlib's Hall marriage
theorem, with Hall's condition derived by a weighted double count, so no integral-flow theorem is
needed. A converse shows that $k\mid n$ is necessary as soon as one colour class is a perfect
matching. The library has 617 lines in three files, contains no \lean{sorry}, and its 32 checked
declarations use only Lean's three standard axioms; with the two verification files added in
version~2 (restatements in Mathlib's own vocabulary and a comparison with independently written
textbook statements) the gate checks 76 declarations, again with standard axioms only. As far as we could determine, within the search
scope described in Section~\ref{sec:related}, no earlier formal proof of the theorem was public.
\end{abstract}

\section{Introduction}

Let $n$ and $k$ be positive integers. The \emph{complete $k$-uniform hypergraph} $K_n^{(k)}$ has $n$
vertices, and every $k$-element set of vertices is an edge. A \emph{1-factor} (perfect matching,
or \emph{parallel class}) of $K_n^{(k)}$ is a set of pairwise disjoint edges that covers every
vertex; it has exactly $n/k$ edges, so it can exist only when $k\mid n$. A \emph{1-factorisation}
is a partition of all $\binom nk$ edges into 1-factors. Counting edges, it has
$\binom nk/(n/k)=\binom{n-1}{k-1}$ classes.

\begin{theorem}[Baranyai~\cite{Baranyai1975}]\label{thm:baranyai}
If $k\mid n$, then $K_n^{(k)}$ has a 1-factorisation: the $k$-subsets of an $n$-set can be
partitioned into $\binom{n-1}{k-1}$ parallel classes.
\end{theorem}

For $k=2$ the theorem says that $K_n$ has a 1-factorisation when $n$ is even, which has been known
since the nineteenth century and is the basis of round-robin tournament schedules: with $n$ teams,
the $n-1$ classes are the rounds, and in each round every team plays exactly one
match~\cite{MendelsohnRosa}. The case $k=3$ (``tournaments of games for three players'') was
settled by Peltesohn in 1936~\cite{Peltesohn}. The general case was open until Baranyai proved it
in 1973; the proof appeared in the Keszthely colloquium volume of 1975~\cite{Baranyai1975}, and
Baranyai developed the method further in~\cite{Baranyai1979}. The textbook proof, due to Brouwer and
Schrijver~\cite{BrouwerSchrijver}, is reproduced by van Lint and Wilson~\cite[Ch.~38]{vLW}. It is
an induction on the vertices in which each step is an integral rounding of a fractional
assignment, obtained there from the integrality theorem for network flows. In design-theoretic
language, the theorem says that the complete design of all $k$-subsets of an $n$-set is
resolvable whenever $k\mid n$.

The theorem is a textbook example of fractional relaxation followed by integral
rounding~\cite[Ch.~38]{vLW}.
Mathlib's list of ``1000+ theorems'' contains it: the file \path{docs/1000.yaml} at the
revision used here (\texttt{0df444a3}) has the entry \texttt{Q790236} with the single field
\texttt{title: Baranyai's theorem} and no declaration or author recorded~\cite{Mathlib1000}, and the
corresponding page of the 1000+ theorems project lists no formalisation~\cite{Plus1000}. It is not
in Mathlib's list of 100 theorems (\path{docs/100.yaml}).

\paragraph{Contributions.}
\begin{itemize}
\item A Lean~4 proof of Baranyai's theorem for the vertex type \lean{Fin n}
(Theorem~\ref{thm:formal}), with the number of colours written as $\binom{n-1}{k-1}$ and, in a
variant, as $\binom nk k/n$; and its specialisation to 1-factorisations of $K_n$.
\item Transport of the theorem to an arbitrary finite type of cardinality $n$ and to the elements
of a \lean{Finset} (Section~\ref{sec:transport}).
\item A partition form (Theorem~\ref{thm:partition}) that states the conclusion in terms of colour
classes: distinct classes are disjoint, the classes cover all $k$-subsets, blocks of one class are
pairwise disjoint, every class has $n/k$ blocks, and the number of classes is $\binom nk/(n/k)$
(for $n>0$, $k>0$, $k\mid n$); and the converse that $k\mid n$ is necessary.
\item A verification package: an axiom gate, six tampering tests, kernel-checked concrete
factorisations, and a statement-level audit (Section~\ref{sec:verify}); and, in version~2, the main
theorems restated with Mathlib-only types (finite partitions, perfect matchings of the complete
\lean{SimpleGraph}) with explicit witnesses, and a comparison with textbook statements formalised by a
writer that saw neither code nor paper (Section~\ref{sec:l1415}).
\end{itemize}

\paragraph{Availability.}
The code is at \url{\repo}. Line references in this paper are to commit \texttt{\shortcommit}
(\texttt{\commit}). The toolchain is Lean \texttt{v4.33.1}~\cite{Lean4} with Mathlib~\cite{Mathlib}
revision \texttt{0df444a3} (tag \texttt{v4.33.1}), pinned in \path{lake-manifest.json}. The library
\lean{Baranyai} consists of \path{Baranyai/Basic.lean} (431 lines), \path{Baranyai/Transport.lean}
(56 lines) and \path{Baranyai/Partition.lean} (130 lines), with a three-line root file. The four
main theorems of \path{Basic.lean} are declared at the top level; all other declarations of these
three files are in the namespace \lean{Baranyai}. Version~2 of this paper adds two verification
files, \path{Baranyai/L14.lean} (409 lines, top-level declarations) and
\path{Baranyai/L15Classical.lean} (358 lines, namespace \lean{Baranyai.Classical}), imported by two
new lines at the end of the root file, in a new commit \texttt{1ab21623da9572cb10a8630a086c010b9e554269}
(Section~\ref{sec:l1415}). The three library files and every statement of commit
\texttt{\shortcommit} are unchanged in that commit, so all line references below remain valid. Table~\ref{tab:thms} lists the original development's informal statements with
their formal counterparts; Section~\ref{sec:l1415} identifies the additional verification statements.
Supporting material (the Layer-14 and Layer-15 reports, the blind writer's statements and report, the transcript audit, and the reviews of this version with the fixes made in response) is in the folder \texttt{Baranyai} of the repository \url{https://github.com/jpbald93/LEAN-verification}.

{\small
\setlength{\LTcapwidth}{\textwidth}
\begin{longtable}{@{}>{\raggedright\arraybackslash}p{0.42\textwidth}
  >{\raggedright\arraybackslash}p{0.38\textwidth}l@{}}
\caption{Informal statements and their formal counterparts. Line numbers refer to the first line
of the declaration at commit \texttt{\shortcommit}.}\label{tab:thms}\\
\toprule
Informal statement & Lean name & Source \\ \midrule
\endfirsthead
\toprule
Informal statement & Lean name & Source \\ \midrule
\endhead
\bottomrule
\endfoot
\multicolumn{3}{@{}l}{\emph{Main theorem (Section~\ref{sec:statement})}}\\
Baranyai's theorem on $\{0,\dots,n-1\}$ (Thm.~\ref{thm:formal}) & \decl{baranyai} & \src{Basic}{396} \\
The same with $\binom nk k/n$ colours, $n>0$ & \decl{baranyai'} & \src{Basic}{404} \\
1-factorisation of $K_n$, $n$ even & \decl{complete_graph_one_factorisation} & \src{Basic}{419} \\
$K_4$ splits into 3 perfect matchings & \decl{baranyai_four_two} & \src{Basic}{428} \\
\multicolumn{3}{@{}l}{\emph{Proof (Section~\ref{sec:proof})}}\\
Stage-$j$ invariant (Def.~\ref{def:inv}) & \decl{Baranyai.Inv} & \src{Basic}{25} \\
Demand $d_j(S)=\binom{n-j-1}{k-|S|-1}$ for $|S|<k$, else $0$ & \decl{Baranyai.dem} & \src{Basic}{33} \\
$(t+1)\binom{N+1}{t+1}=(N+1)\binom Nt$ & \decl{Baranyai.choose_key} & \src{Basic}{36} \\
$(k-|S|)\binom{n-j}{k-|S|}=(n-j)\,d_j(S)$ (Lemma~\ref{lem:arith}) & \decl{Baranyai.weight_eq} & \src{Basic}{40} \\
$\binom{n-j}{k-|S|}=d_j(S)+\binom{n-j-1}{k-|S|}$ (Lemma~\ref{lem:arith}) & \decl{Baranyai.pascal_dem} & \src{Basic}{52} \\
$n\binom{n-1}{k-1}=\binom nk k$ & \decl{Baranyai.choose_mul} & \src{Basic}{64} \\
The slots of a class have total size $j$ & \decl{Baranyai.row_card_sum} & \src{Basic}{71} \\
$\sum_p (k-|A_{ip}|)=n-j$ & \decl{Baranyai.row_weight_sum} & \src{Basic}{78} \\
Weighted double count (Lemma~\ref{lem:wc}) & \decl{Baranyai.weighted_count} & \src{Basic}{84} \\
Induction step, via Hall (Lemma~\ref{lem:step}) & \decl{Baranyai.step} & \src{Basic}{106} \\
Stage $0$: all slots empty & \decl{Baranyai.init} & \src{Basic}{304} \\
Stage $n$ yields the factorisation (Lemma~\ref{lem:final}) & \decl{Baranyai.exists_partition_of_inv} & \src{Basic}{325} \\
An invariant exists at stage $n$ & \decl{Baranyai.exists_final} & \src{Basic}{374} \\
\multicolumn{3}{@{}l}{\emph{Transport (Section~\ref{sec:transport})}}\\
Factorisations transport along $\alpha\simeq\beta$ & \decl{Baranyai.factorisation_congr} & \src{Transport}{17} \\
Baranyai for any finite type with $|\alpha|=n$ & \decl{Baranyai.baranyai_of_card} & \src{Transport}{43} \\
Baranyai for the elements of a \lean{Finset} & \decl{Baranyai.baranyai_finset} & \src{Transport}{50} \\
\multicolumn{3}{@{}l}{\emph{Partition form (Section~\ref{sec:partition})}}\\
Distinct colour classes are disjoint & \decl{Baranyai.class_disjoint} & \src{Partition}{23} \\
The colour classes cover all $k$-subsets & \decl{Baranyai.classes_cover} & \src{Partition}{31} \\
Two blocks of one colour are disjoint & \decl{Baranyai.block_disjoint} & \src{Partition}{38} \\
$(\text{blocks of a class})\cdot k=|\alpha|$ & \decl{Baranyai.class_card_mul} & \src{Partition}{48} \\
A class has $|\alpha|/k$ blocks, for $k>0$ & \decl{Baranyai.class_card} & \src{Partition}{71} \\
$\binom{n-1}{k-1}=\binom nk/(n/k)$, for $n>0$, $k>0$, $k\mid n$ & \decl{Baranyai.num_classes} & \src{Partition}{78} \\
Partition form of Baranyai (Thm.~\ref{thm:partition}) & \decl{Baranyai.baranyai_partition} & \src{Partition}{89} \\
Converse: one perfect-matching class forces $k\mid|\alpha|$ & \decl{Baranyai.dvd_card_of_factorisation} & \src{Partition}{106} \\
\end{longtable}
}

\section{The theorem and its formal statement}\label{sec:statement}

Vertices are the elements of \lean{Fin n}, that is $\{0,\dots,n-1\}$. A $k$-subset is an element of
the subtype \lean{\{S : Finset (Fin n) // S.card = k\}}, and a candidate factorisation is a
colouring \lean{c} of the $k$-subsets with colours in \lean{Fin M}. The factorisation property is
stated directly: for every colour $i$ and every vertex $v$ there is exactly one $k$-subset of colour
$i$ that contains $v$. The main theorem is
\src{Basic}{396}:
\begin{lstlisting}
theorem baranyai (n k : ℕ) (hk : 0 < k) (hkn : k ∣ n) :
    ∃ c : {S : Finset (Fin n) // S.card = k} → Fin ((n - 1).choose (k - 1)),
      ∀ i (v : Fin n), ∃! S : {S : Finset (Fin n) // S.card = k}, c S = i ∧ v ∈ S.1
\end{lstlisting}

\begin{theorem}[\decl{baranyai}]\label{thm:formal}
Let $0<k$ and $k\mid n$. There is a map $c$ from the $k$-subsets of $\{0,\dots,n-1\}$ to
$\{0,\dots,\binom{n-1}{k-1}-1\}$ such that for every colour $i$ and every vertex $v$ exactly one
$k$-subset $S$ has $c(S)=i$ and $v\in S$.
\end{theorem}

\paragraph{Reading the statement.} Let $n>0$. The map $c$ assigns each $k$-subset exactly one
colour, and every colour class $c^{-1}(i)$ is non-empty, since it contains a block through the
vertex $0$; so the classes form a partition of the set of all $k$-subsets into
$\binom{n-1}{k-1}$ (non-empty) parts. The condition ``exactly one block of colour $i$ through each
vertex'' says precisely that the class $c^{-1}(i)$ is a set of pairwise disjoint $k$-subsets
covering every vertex, i.e.\ a 1-factor. Hence, for $n>0$, the conclusion is that $K_n^{(k)}$ has a
1-factorisation with $\binom{n-1}{k-1}$ classes, as in Theorem~\ref{thm:baranyai}. The extreme
cases agree with the code: for $k=1$ there is $\binom{n-1}0=1$ colour, whose class is the set of
all $n$ singletons, and for $k=n$ there is $\binom{n-1}{n-1}=1$ colour, whose class is the single
block $\{0,\dots,n-1\}$. Section~\ref{sec:partition} proves the parts of this reading that are not
literally in the statement (disjointness of blocks, the size of each class). Natural-number
subtraction is truncated in Lean; for $n,k\ge1$, $n-1$ and $k-1$ have their usual meaning.

The case $n=0$ must be read separately. There are no vertices and no $k$-subsets, and
\lean{(0 - 1).choose (k - 1)} is $1$ for $k=1$ and $0$ for $k\ge2$. For $k=1$ the code therefore
supplies one colour whose fibre is empty; since the parts of an ordinary partition are non-empty,
this family of fibres is not a partition with the advertised number of parts. What
\decl{baranyai} proves for $n=0$ is only the valid but vacuous colouring statement: the (empty)
map $c$ exists and the condition ``for every colour $i$ and every vertex $v$\,\dots'' holds because
there is no vertex. \decl{baranyai'}, which writes the number of colours as
\lean{n.choose k * k / n}, and the partition form (Theorem~\ref{thm:partition}) assume $0<n$.

The hypothesis $0<k$ is not logically required: for $k=0$ the hypothesis $k\mid n$ forces $n=0$,
and the statement holds vacuously (there is no vertex). It is used by the proof, and with $k=0$
the theorem would say nothing. The hypothesis $k\mid n$ is necessary, since a single 1-factor
partitions the $n$ vertices into $k$-sets. This is a theorem of the library,
\decl{Baranyai.dvd_card_of_factorisation} \src{Partition}{106}: if some colour class of a colouring
of the $k$-subsets of a finite type $\alpha$ covers every vertex exactly once, then $k\mid|\alpha|$.
The tampering test \lean{falsehyp} (Section~\ref{sec:verify}) deletes \lean{hkn} from the statement
and checks that the proof no longer compiles; this shows only that the proof uses \lean{hkn}. That
the weakened statement is false is a separate fact: for $n=3$, $k=2$ three vertices cannot be
covered by disjoint pairs, and \path{scratch/Referee.lean} checks with \lean{decide} that no
colouring of the 2-subsets of \lean{Fin 3} with two colours has the factorisation property.

For $k=2$ the theorem specialises to \decl{complete_graph_one_factorisation} \src{Basic}{419}:
if $n$ is even there is a colouring of the 2-subsets of \lean{Fin n} with $n-1$ colours in which
every colour class is a perfect matching. \decl{baranyai_four_two} is the instance $n=4$.

\paragraph{Transport and partition statements.} \path{Transport.lean} replaces \lean{Fin n} by an
arbitrary finite type \src{Transport}{43}:
\begin{lstlisting}
theorem baranyai_of_card [Fintype α] (n k : ℕ) (hk : 0 < k) (hkn : k ∣ n)
    (hα : Fintype.card α = n) :
    ∃ c : {S : Finset α // S.card = k} → Fin ((n - 1).choose (k - 1)),
      ∀ i (v : α), ∃! S : {S : Finset α // S.card = k}, c S = i ∧ v ∈ S.1
\end{lstlisting}
and \path{Partition.lean} states the conclusion as a partition \src{Partition}{89}:
\begin{lstlisting}
theorem baranyai_partition [Fintype α] [DecidableEq α] (n k : ℕ) (hk : 0 < k) (hn : 0 < n)
    (hkn : k ∣ n) (hα : Fintype.card α = n) :
    ∃ c : {S : Finset α // S.card = k} → Fin ((n - 1).choose (k - 1)),
      (∀ i (v : α), ∃! S : {S : Finset α // S.card = k}, c S = i ∧ v ∈ S.1) ∧
      (∀ S T, S ≠ T → c S = c T → Disjoint S.1 T.1) ∧
      (∀ i j, i ≠ j →
        Disjoint (univ.filter fun S => c S = i) (univ.filter fun S => c S = j)) ∧
      univ.biUnion (fun i => univ.filter fun S => c S = i) = univ ∧
      (∀ i, (univ.filter fun S => c S = i).card = n / k) ∧
      (n - 1).choose (k - 1) = n.choose k / (n / k)
\end{lstlisting}

\section{Proof outline}\label{sec:proof}

We follow the proof of Brouwer and Schrijver~\cite{BrouwerSchrijver} as given in~\cite[Ch.~38]{vLW},
with one change in the rounding step. Write $n=mk$ and $M=\binom{n-1}{k-1}$, and let
$[j]=\{0,\dots,j-1\}$. The formal \decl{Baranyai.Inv}, \decl{Baranyai.step} and
\decl{Baranyai.exists_partition_of_inv} are stated for arbitrary $M$ and $m$ under their own
hypotheses (for \decl{Baranyai.step} only the stage-$j$ invariant \lean{hI : Inv n k m M j A}, $j<n$
and $n=mk$; for
\decl{Baranyai.exists_partition_of_inv} only the stage-$n$ invariant), and
\decl{Baranyai.block_disjoint} needs no finite-type assumption on the vertices, so the argument
below and in Section~\ref{sec:partition} specialises these more general results. Binomial coefficients with a negative lower index do not occur: all lower
indices below are natural numbers, and $\binom ab=0$ for $b>a$.

\begin{definition}[\decl{Baranyai.Inv}, \src{Basic}{25}]\label{def:inv}
A family $A=(A_{ip})$, $i<M$, $p<m$, of subsets of $\{0,\dots,n-1\}$ satisfies the
\emph{stage-$j$ invariant} if
\begin{enumerate}
\item for each class $i$ the slots $A_{i0},\dots,A_{i,m-1}$ are pairwise disjoint (\lean{disj});
\item for each class $i$ their union is $[j]$ (\lean{cover});
\item every slot has at most $k$ elements (\lean{size});
\item every $S\subseteq[j]$ with $|S|\le k$ occurs as a slot exactly $\binom{n-j}{k-|S|}$ times,
i.e.\ $\#\{(i,p):A_{ip}=S\}=\binom{n-j}{k-|S|}$ (\lean{cnt}).
\end{enumerate}
\end{definition}
Slots may be empty, and an empty slot is counted once per index $(i,p)$; this is how the formal
version realises the multisets of~\cite{vLW}. At stage $0$ all slots are empty, and the
count condition for $S=\varnothing$ reads $Mm=\binom nk$, which follows from
$n\binom{n-1}{k-1}=\binom nk k$ (\decl{Baranyai.choose_mul}); this is \decl{Baranyai.init}. At stage
$n$ the theorem can be read off:

\begin{lemma}[\decl{Baranyai.exists_partition_of_inv}]\label{lem:final}
If $A$ satisfies the stage-$n$ invariant, then a factorisation with $M$ colours exists: colour each
$k$-subset $S$ by the class $i$ of its slot. (The formal statement asserts only existence.)
\end{lemma}
\begin{proof}[Sketch]
At $j=n$ condition~4 says that $S$ occurs $\binom{0}{k-|S|}$ times: once if $|S|=k$ and never if
$|S|<k$. Every slot occurs at least once, so every slot has exactly $k$ elements, and every
$k$-subset is exactly one slot $(i,p)$. Colour $S$ by $i$. Given a colour $i$ and a vertex $v$,
condition~2 gives a slot of class $i$ containing $v$, and condition~1 makes it unique.
\end{proof}

\paragraph{The induction step.} Let $j<n$. Each class must choose one of its slots to receive the
new vertex $j$. Define the \emph{demand} $d_j(S)=\binom{n-j-1}{k-|S|-1}$ if $|S|<k$ and
$d_j(S)=0$ otherwise (\decl{Baranyai.dem}). The step works if each set $S$ is chosen exactly
$d_j(S)$ times: then, for $j\in S$, the new count of $S$ is the number of times $S\setminus\{j\}$
was chosen, $d_j(S\setminus\{j\})=\binom{n-j-1}{k-|S|}$, and for $j\notin S$ it is the old count
minus the number of times $S$ was chosen, which is again $\binom{n-j-1}{k-|S|}$ by the next lemma.

\begin{lemma}[\decl{Baranyai.weight_eq}, \decl{Baranyai.pascal_dem}]\label{lem:arith}
For $j<n$ and every set $S$ with $|S|\le k$,
\[(k-|S|)\binom{n-j}{k-|S|}=(n-j)\,d_j(S),\qquad
\binom{n-j}{k-|S|}=d_j(S)+\binom{n-j-1}{k-|S|}.\]
\end{lemma}
The first identity is the absorption identity $(t+1)\binom{N+1}{t+1}=(N+1)\binom Nt$
(\decl{Baranyai.choose_key}, from Mathlib's \mdecl{Nat.add_one_mul_choose_eq}); the second is
Pascal's rule (\mdecl{Nat.choose_succ_succ'}). In both cases $|S|=k$ is handled separately. The
formal \decl{Baranyai.weight_eq} has no size hypothesis at all, and the hypothesis $|S|\le k$ of
\decl{Baranyai.pascal_dem} turns out not to be used by its proof (Section~\ref{sec:notes}).

\paragraph{Weighted double counting.} In~\cite{vLW} the choice is made by exhibiting a fractional
flow, in which class $i$ sends weight $(k-|S|)/(n-j)$ to each of its slots $S$, and then appealing
to the integrality of maximum flows (their Theorem~7.2, or alternatively their rounding theorem
for circulations, Theorem~7.4). The formal proof uses the same
weights, cleared of denominators. Since the slots of class $i$ partition $[j]$ and there are $m$
of them, $\sum_p(k-|A_{ip}|)=mk-j=n-j$ (\decl{Baranyai.row_weight_sum}, via
\decl{Baranyai.row_card_sum}).

\begin{lemma}[\decl{Baranyai.weighted_count}]\label{lem:wc}
Assume the stage-$j$ invariant, $j<n$, and let $\mathcal Y$ be a family of subsets of $[j]$ of size
at most $k$. Then
\[\sum_{(i,p):\,A_{ip}\in\mathcal Y}(k-|A_{ip}|)=(n-j)\sum_{S\in\mathcal Y}d_j(S).\]
\end{lemma}
\begin{proof}
Group the left-hand side by the value $S=A_{ip}$ (\mdecl{Finset.sum_fiberwise_of_maps_to}). The
fibre of $S$ has $\binom{n-j}{k-|S|}$ elements by condition~4, each contributing $k-|S|$; apply the
first identity of Lemma~\ref{lem:arith}.
\end{proof}

\paragraph{Selection by Hall's theorem.} For each class $i$ let $t(i)$ be the finite set of
\emph{tokens} $(S,r)$ with $S$ a slot of class $i$ and $r<d_j(S)$; in Lean this is
\lean{(univ.image (A i)).sigma (fun S => range (dem n k j S))}. A system of distinct
representatives for $t$ is exactly a choice of one slot per class in which each $S$ is chosen at
most $d_j(S)$ times. Mathlib's Hall theorem, in the form
\mdecl{Finset.all_card_le_biUnion_card_iff_exists_injective}~\cite{Hall,HallLean}, says that such a
system exists iff $|X|\le|\bigcup_{i\in X}t(i)|$ for every set $X$ of classes.

\begin{lemma}[\decl{Baranyai.step}]\label{lem:step}
If $A$ satisfies the stage-$j$ invariant, $j<n$ and $n=mk$, then some $A'$ satisfies the
stage-$(j+1)$ invariant.
\end{lemma}
\begin{proof}[Sketch]
\emph{Hall's condition.} Let $X$ be a set of classes and $\mathcal Y$ the family of sets occurring
as slots of classes in $X$. Then $\bigcup_{i\in X}t(i)$ is the set of tokens $(S,r)$ with
$S\in\mathcal Y$, $r<d_j(S)$, of size $\sum_{S\in\mathcal Y}d_j(S)$. By \decl{Baranyai.row_weight_sum}
and Lemma~\ref{lem:wc},
\[(n-j)|X|=\sum_{i\in X}\sum_p(k-|A_{ip}|)\le\sum_{(i,p):\,A_{ip}\in\mathcal Y}(k-|A_{ip}|)
=(n-j)\sum_{S\in\mathcal Y}d_j(S),\]
and $n-j>0$. \emph{Exactness.} Hall gives an injective choice $i\mapsto(S_i,r_i)\in t(i)$; let
$p(i)$ be a slot of class $i$ with $A_{i,p(i)}=S_i$. Injectivity in $r$ bounds the number of
classes choosing $S$ by $d_j(S)$ (\mdecl{Finset.card_le_card_of_injOn}). Summed over the family
$\mathcal U$ of all subsets of $[j]$ of size at most $k$, the numbers of choices add up to $M$
(\mdecl{Finset.card_eq_sum_card_fiberwise}). Lemma~\ref{lem:wc} with $\mathcal Y=\mathcal U$, whose
left side is $M(n-j)$, gives $\sum_{S\in\mathcal U}d_j(S)=M$. A sum of termwise inequalities with
equal totals is a sum of equalities (\mdecl{Finset.sum_eq_sum_iff_of_le}), so every $S$ is chosen
exactly $d_j(S)$ times. A chosen slot has $d_j(S)>0$ and hence $|S|<k$. \emph{Update.} Put
$A'_{i,p(i)}=A_{i,p(i)}\cup\{j\}$ and leave the other slots unchanged. Conditions~1--3 are
immediate; condition~4 is the computation before Lemma~\ref{lem:arith}, split according to whether
$j\in S$.
\end{proof}
\decl{Baranyai.exists_final} iterates Lemma~\ref{lem:step} from \decl{Baranyai.init} up to $j=n$, and
\decl{baranyai} combines it with Lemma~\ref{lem:final}.

\paragraph{Hall instead of flows.} Because every class has capacity one and the demand of each set
can be split into $d_j(S)$ unit tokens, the integral rounding needed here is a bipartite matching
problem, and Hall's theorem suffices. Hall's theorem is itself a consequence of flow integrality,
and replacing a capacitated assignment by a matching through token copies is a standard device; we
do not claim either as new. Its use here means that no theory of network flows had to be
formalised: the integrality of the solution comes from Hall's theorem, and the fractional flow
survives only in the double count of Lemma~\ref{lem:wc}.

\paragraph{What was reused and what was built.} From Mathlib we use, among others, Hall's marriage theorem
(\mdecl{Finset.all_card_le_biUnion_card_iff_exists_injective}, formalised by Gusakov, Mehta and
Miller~\cite{HallLean}), the binomial identities \mdecl{Nat.add_one_mul_choose_eq} and
\mdecl{Nat.choose_succ_succ'}, the \lean{Finset} big-operator lemmas named above together with
\mdecl{Finset.card_biUnion} and \mdecl{Finset.sum_tsub_distrib}, and, for the transport, the
equivalences \mdecl{Equiv.finsetCongr} and \mdecl{Fintype.equivFinOfCardEq} and the lemma
\mdecl{Fintype.card_coe}; the proofs also use standard \lean{Finset} and \lean{Nat} lemmas such as
\lean{card\_image\_of\_injective}, \lean{card\_union\_of\_disjoint}, \lean{sum\_le\_sum\_of\_subset},
\lean{sum\_product}, \lean{Fintype.sum\_prod\_type}, \lean{Fin.card\_filter\_val\_lt},
\lean{card\_erase\_add\_one} and \lean{Nat.choose\_eq\_zero\_of\_lt}. Everything else was written for this development: the invariant and the
demand function, the arithmetic and double-counting lemmas, the induction step and its
initialisation, the extraction of the factorisation, the transport and the partition corollaries.

\section{Transport and partition form}\label{sec:new}

\subsection{Transport}\label{sec:transport}
\decl{Baranyai.factorisation_congr} \src{Transport}{17} says that if the $k$-subsets of $\alpha$ have
a factorisation with $M$ colours, so do the $k$-subsets of $\beta$, for any bijection
$e:\alpha\simeq\beta$. The colour of $T\subseteq\beta$ is the colour of its preimage, computed with
\mdecl{Equiv.finsetCongr}, which preserves cardinality and transports membership.
\decl{Baranyai.baranyai_of_card} \src{Transport}{43} applies this to
\mdecl{Fintype.equivFinOfCardEq} and gives Theorem~\ref{thm:formal} for any finite type of
cardinality $n$. \decl{Baranyai.baranyai_finset} \src{Transport}{50} is the case of the subtype
$\uparrow s$ of a \lean{Finset} $s$, with $|s|$ in place of $n$; the $k$-subsets are then
\lean{Finset}s of elements of $s$.

These statements add no mathematics: Theorem~\ref{thm:baranyai} is about an abstract $n$-set, and
the transport removes the artefact of the vertex type \lean{Fin n}. They make the theorem usable
for vertex sets that are not literally $\{0,\dots,n-1\}$; the satisfiability file instantiates them
for \lean{\{1, 5, 7, 9\} : Finset}~$\NN$ and for the type \lean{Fin 2 × Fin 3}.

\subsection{Partition form}\label{sec:partition}
The lemmas of \path{Partition.lean} hold for an arbitrary finite type $\alpha$ (for
\decl{Baranyai.block_disjoint}, an arbitrary type) and any colouring $c$ of its $k$-subsets with
colours in \lean{Fin M}.
\begin{itemize}
\item \decl{Baranyai.class_disjoint}: distinct colour classes $c^{-1}(i)$, $c^{-1}(j)$ are disjoint.
\item \decl{Baranyai.classes_cover}: the union of all colour classes is the set of all $k$-subsets.
\end{itemize}
These two hold for any function $c$: the fibres are pairwise disjoint and cover the set of
$k$-subsets, and the non-empty fibres form a partition of it (a colour that is not used has an
empty fibre, which can happen for an arbitrary $c$ even when $n>0$). The next three use the factorisation property \lean{hc}.
\begin{itemize}
\item \decl{Baranyai.block_disjoint}: two different $k$-subsets of the same colour are disjoint.
\item \decl{Baranyai.class_card_mul}: every class has $b$ blocks with $b\cdot k=|\alpha|$, by double
counting incidences between the blocks of the class and the vertices, each vertex lying in
exactly one block. \decl{Baranyai.class_card} divides by $k$ and therefore assumes $k>0$.
\end{itemize}
Finally \decl{Baranyai.num_classes} proves $\binom{n-1}{k-1}=\binom nk/(n/k)$ under the
hypotheses $n>0$, $k>0$ and $k\mid n$, the edge count of the introduction. The hypothesis $n>0$ is
needed: for $n=0$, $k=1$ the left side is $\binom{0}{0}=1$ and the right side is
$\binom01/(0/1)=0$. \decl{Baranyai.baranyai_partition} (Theorem~\ref{thm:partition}) collects all of this into one statement.
The same double count gives the converse \decl{Baranyai.dvd_card_of_factorisation}: if one colour
class covers every vertex exactly once, then $k$ divides $|\alpha|$.

\begin{theorem}[\decl{Baranyai.baranyai_partition}]\label{thm:partition}
Let $\alpha$ be a finite type with $n$ elements, $0<k$, $0<n$ and $k\mid n$. There is a colouring
$c$ of the $k$-subsets of $\alpha$ with $\binom{n-1}{k-1}$ colours such that every vertex lies in
exactly one block of each colour; two different blocks of the same colour are disjoint; distinct
colour classes are disjoint and together contain all $k$-subsets; every colour class has $n/k$
blocks; and $\binom{n-1}{k-1}=\binom nk/(n/k)$.
\end{theorem}

The first conjunct is the conclusion of Theorem~\ref{thm:formal}; the others follow from it, so
Theorem~\ref{thm:partition} is not a stronger theorem. Its purpose is to make the informal reading
of Section~\ref{sec:statement} checkable: a reader who prefers to see ``each class is a set of
$n/k$ pairwise disjoint blocks, and the classes partition the edge set'' finds exactly that in
the statement. As the header of \path{Partition.lean} says, the hard direction is not reproved.

\section{Formalisation notes}\label{sec:notes}

\paragraph{Representation.} Vertices are \lean{Fin n} in the core development. Edges are
\lean{Finset}s with a cardinality constraint, packaged as a subtype, rather than \lean{Set}s; this
makes the families of edges finite types and lets \lean{decide} check small instances. A
factorisation is a function from edges to colours, so the property ``each edge lies in exactly
one class'' holds by construction and only the 1-factor property has to be stated. In the proof,
classes and slots are indexed by \lean{Fin M} and \lean{Fin m}, and the invariant is a structure
with four \lean{Prop} fields (\decl{Baranyai.Inv}). Counting slots by index rather than using
multisets removes the need to count the empty set with multiplicity. Hall's theorem is applied to
the token sets $t(i)$ of type \lean{Finset ((\_ : Finset (Fin n)) }$\times$\lean{ }$\mathbb{N}$\lean{)}, a dependent pair of a set
and a token number, so that the injective function produced by Hall encodes both the chosen slot
and the multiplicity bound.

\paragraph{Bridges to Mathlib's graph and hypergraph notions.} The factorisation property is a
direct $\exists!$ statement about \lean{Finset}s, and the three library files do not connect it to
Mathlib's \mdecl{SimpleGraph.Subgraph.IsPerfectMatching}, to the structure \mdecl{Hypergraph} (which
exists at the pinned revision, with \lean{Set}-valued edges), or to \mdecl{Finpartition}. In
particular \decl{complete_graph_one_factorisation} is stated in the same $\exists!$ form. The
verification files of version~2 supply the bridges to \mdecl{Finpartition} and to perfect matchings:
\path{L14.lean} restates the main theorems as families of \mdecl{Finpartition}s and the case $k=2$ as
a decomposition of the complete \lean{SimpleGraph} into subgraphs satisfying
\mdecl{SimpleGraph.Subgraph.IsPerfectMatching}, and \path{L15Classical.lean} proves the textbook forms
stated with these notions (Section~\ref{sec:l1415}). No bridge to \mdecl{Hypergraph} has been
formalised. The proof devices \decl{Baranyai.Inv} and
\decl{Baranyai.dem} are local and are not bridged to anything; they do not occur in the statements
of the main theorems.

\paragraph{An unused hypothesis.} The statement-level audit (Section~\ref{sec:verify}) found that the
hypothesis \lean{hS : S.card }$\le$\lean{ k} of \decl{Baranyai.pascal_dem} is not used by its proof: the lemma
holds without it (when $|S|\ge k$ both sides reduce to $\binom{n-j}0=1=0+1$ under truncated
subtraction).

\paragraph{Size.} \path{Basic.lean} has 431 lines, of which the induction step takes about 200. The
two corollary files have 186 lines together. In these three files the definitions are
\decl{Baranyai.Inv} and \decl{Baranyai.dem}; all other declarations are theorems. (The Layer-15 file
of version~2 adds the writer's four definitions, two bridge definitions, and a proposition-valued
definition recording the unproved general almost-regular statement.)

\section{Verification}\label{sec:verify}

Except for the version-2 checks explicitly identified below, the checks in this section were run on
commit \texttt{\shortcommit}. The library sources, \path{gate.sh},
\path{tests/tamper.sh}, \path{scratch/Satisfiable.lean} and \path{scratch/Referee.lean} on the build machine were confirmed
byte-identical (SHA-256) to those of the commit before the runs.

\paragraph{Axiom gate.} \file{gate.sh} rejects the tokens \lean{sorry}, \lean{admit},
\lean{native\_decide}, \lean{axiom}, \lean{\#eval}, \lean{run\_cmd}, \lean{set\_option}, \lean{macro},
\lean{elab}, \lean{syntax}, \lean{notation} and \lean{import Lean} in the library sources, builds the
library (8\,710 jobs, no errors), generates its own \lean{\#print axioms} query for each of 32 fully
qualified declarations (every theorem, definition and structure field of the library), and fails on
any axiom other than \lean{propext}, \lean{Classical.choice} and \lean{Quot.sound}. It prints
\texttt{PASS (32 declarations, standard axioms only)}. With the two files of version~2 it prints
\texttt{PASS (76 declarations, standard axioms only)}, and the tampering tests below are again all
rejected (Section~\ref{sec:l1415}).

\paragraph{Tampering tests.} \file{tests/tamper.sh} plants six defects in scratch copies of the
repository and requires the gate to fail on each with a nonzero exit status and the expected message:
a \lean{sorry}; an \lean{axiom} split over two lines; an \lean{\#eval} printing a fake axiom report;
a macro rewriting \lean{\#print axioms}; the statement of \decl{baranyai} with the hypothesis
\lean{k }$\mid$\lean{ n} deleted, which must fail to build (this shows that the proof uses the
hypothesis; that the weakened statement is false is checked separately in \path{scratch/Referee.lean}); and a genuine extra
axiom with a long name, so that Lean wraps the report over several lines. All six are rejected
(6/6).

\paragraph{Non-triviality and instances.} The hypotheses $0<k$, $k\mid n$ of
Theorem~\ref{thm:formal} are visibly satisfiable, so the real risk is a conclusion that is too weak
to fail. Two kernel-checked facts address it. First, \file{scratch/Referee.lean} (outside the
library, exit status~$0$) proves by \lean{decide} that the conclusion fails at $n=3$, $k=2$: no map
from the 2-subsets of \lean{Fin 3} to \lean{Fin 2} has the factorisation property. Two colours
are the count prescribed by Theorem~\ref{thm:formal}, $\binom{2}{1}=2$. Second, the
library corollary \decl{Baranyai.dvd_card_of_factorisation} shows that the conclusion with at least
one colour forces $k\mid n$; \path{scratch/Referee.lean} applies it to rule out every positive colour
count at $n=3$, $k=2$. The same file gives a term of the hypothesis of \decl{Baranyai.step}, which
carries $j<n$: the stage-$0$ invariant \decl{Baranyai.init} at $n=4$, $k=2$, $m=2$, with $0<4$, fed to
\decl{Baranyai.step}. It also checks that three hypotheses used by the proof are logically
removable (see the audit below). The file \file{scratch/Satisfiable.lean}, also outside the library,
compiles with exit status~$0$. It shows that the conclusion is met in concrete cases, which is a
weaker point than non-triviality. It contains explicit factorisations for $(n,k)=(4,2)$
(the matchings $\{01,23\}$, $\{02,13\}$, $\{03,12\}$) and $(6,3)$ (ten classes, each a triple through
$0$ together with its complement), each checked by \lean{decide} and not derived from the theorem;
the hypotheses $0<3$ and $3\mid6$ of \decl{baranyai}, checked by \lean{decide}; an explicit
stage-$4$ witness of \decl{Baranyai.Inv} for $n=4$, $k=2$, $m=2$, $M=3$, meeting the hypothesis of
\decl{Baranyai.exists_partition_of_inv}; and instances of \decl{Baranyai.baranyai_finset},
\decl{Baranyai.baranyai_of_card} and \decl{Baranyai.baranyai_partition}. The Python script
\file{code/check_baranyai.py} runs the same induction numerically, with a capacitated bipartite
matching in the rounding step, and verifies the resulting factorisations for
$(n,k)\in\{(9,3),(8,2),(10,5),(12,4),(12,3)\}$.

\paragraph{Statement-level audit.} A separate audit tool was run on \path{Baranyai/Basic.lean} (the
same file, identified by its SHA-256). The audit has since been re-run over the complete shipped tree --- all three library files, 26
theorems --- with the same tool: 93 conclusion mutations, all rejected; no contradictory
hypothesis set (24 probed, 2 declarations without hypotheses); every recognised hypothesis
classified as needed by its proof except the unused one identified below and the two hypotheses
of \decl{baranyai'}, whose control copy does not compile in either run; and no unused-variable
warning under the linter, in the three library files or in the three files under \path{scratch/}.
The figures in the three parts below are those of the first run on \path{Basic.lean}, whose
individual mutation diagnostics were adjudicated case by case.
\begin{itemize}
\item \emph{Conclusion mutations.} For each of the 15 theorems of the file, small changes are made
to the conclusion (flip a relation, negate an equality, change a numeral, and similar) while the
proof is kept, and Lean must reject the result. All 55 generated edits failed to compile; none
survived. The rejections are of two kinds. In 29 of them the mutated statement is itself ill-typed
and the first diagnostic identifies an error in the edited statement: for example, twelve of them
render as a malformed application such as \lean{n 1}, one of them as \lean{j 1}, and all seven
mutants of \decl{Baranyai.exists_final} are ill-typed (four malformed applications and three
dependent-domain mismatches). These 29 say nothing about the proofs. Only the
remaining 26 mutants are well-typed statements whose proofs Lean rejects, and only these show
that a proof depends on the exact conclusion.
\item \emph{Contradiction probe.} For each theorem with hypotheses, automation (\lean{omega},
\lean{decide}, \lean{simp\_all}, \lean{linarith} and others) tries to derive \lean{False} from the
hypotheses alone, alongside a control that must compile. No hypothesis set was found
contradictory (13 probed, 2 theorems without hypotheses). This probe is weak: the probed sets
are either simple arithmetic facts or contain \lean{hI : Inv \dots}, which the automation cannot
unfold. The real consistency evidence is the stage-$0$ witness from \decl{Baranyai.init}, fed to
\decl{Baranyai.step} in \path{scratch/Referee.lean}, and the stage-$n$ witness \lean{A42\_inv} in
\path{scratch/Satisfiable.lean}.
\item \emph{Hypothesis ablation.} Selected, syntactically recognised hypotheses are replaced in
turn by \lean{True} in a renamed copy of the declaration, which must then fail to compile. The
selector does not cover every hypothesis: it omits, among others, divisibility hypotheses such as
\lean{hkn : k }$\mid$\lean{ n}, invariant hypotheses such as \lean{hI : Inv \dots}, and parity
hypotheses such as \lean{hn : Even n}. For this file it produced 19 jobs, of which 17 were
executed and 2 skipped: 16 hypotheses were found needed \emph{by this proof} and 1 unused. For
\decl{baranyai} the hypothesis \lean{0 < k} is needed by this proof; for \decl{Baranyai.pascal_dem}
the hypothesis \lean{S.card }$\le$\lean{ k} is not (Section~\ref{sec:notes}). Ablation does not test
whether a statement stays true without the hypothesis. \path{scratch/Referee.lean} kernel-checks
that \decl{baranyai} without \lean{0 < k}, \decl{Baranyai.choose_mul} without \lean{0 < n} and
\decl{Baranyai.weight_eq} without \lean{j < n} are still true; these three hypotheses are used by
the proofs but are logically removable. \decl{Baranyai.pascal_dem} is merely the one case where the
proof does not use its hypothesis. The skipped jobs include the
ablation of \decl{baranyai'}, which was skipped because its control copy did not compile. That
the proof uses the hypothesis \lean{hkn} of \decl{baranyai} is shown instead by the tampering test
above; that it is logically necessary is \decl{Baranyai.dvd_card_of_factorisation}.
\end{itemize}
The audit also attempts satisfiability certificates for the hypotheses of each theorem. That pass
was model-assisted: a local language model, treated as untrusted, proposes concrete values and only
Lean decides. In the first run on \path{Basic.lean} it certified 6 theorems and left 7 unresolved,
including \decl{baranyai} and \decl{baranyai'}; in the re-run over the three files it produced no
further certificate (18 unresolved, 6 of them failed to produce a candidate). The hand-written
\path{scratch/Satisfiable.lean} covers the main theorem.

\paragraph{Limits of these checks.} None of the checks above shows that the formal statement says what
the prose says (Section~\ref{sec:l1415} adds an independent comparison with textbook statements). The 26 rejected well-typed conclusion mutants show that those proofs depend on
the exact conclusions; the other 29 rejections only show that the edits did not typecheck. A failed contradiction probe
only shows that automation found no contradiction. Neither check is a proof of satisfiability or of faithfulness. The gate trusts the pinned toolchain, Mathlib and lake
configuration. The comparison of formal and informal statements was done by hand: the author
checked each statement in Table~\ref{tab:thms} against its informal counterpart.
Section~\ref{sec:statement} and Theorem~\ref{thm:partition} are meant to make that comparison easy
for the reader.

\paragraph{Name check.} The author's script \path{namegate.sh} (not included in the verification repository) checks that every
declaration typeset with \lean{\textbackslash decl} exists in the local source tree at the commit, in
the stated namespace; that every Mathlib name typeset with \lean{\textbackslash mdecl} is in a list of
names checked with \lean{\#check} against the pinned Mathlib; and that every source link lands on the
declaration line of the declaration it accompanies. It passes on commit \texttt{\shortcommit} for
every name outside Section~\ref{sec:l1415}; the 15 newly introduced declaration names cited in that section are declared only in the new
commit. Run against that commit, the script finds the six top-level ones; the nine in the
sub-namespace \lean{Baranyai.Classical}, which its namespace parser does not handle, were confirmed by
a direct search of \path{L15Classical.lean}.

\subsection{Layers 14 and 15}\label{sec:l1415}

The two files of this subsection were added in version~2, in commit \texttt{1ab21623da9572cb10a8630a086c010b9e554269}; the three
original library source files and their statements, together with the tampering script, are unchanged (the hashes of \path{Basic.lean}, \path{Transport.lean},
\path{Partition.lean} and \path{tests/tamper.sh} are those of \texttt{\shortcommit}). On the new commit
the library builds, the gate prints \texttt{PASS (76 declarations, standard axioms only)} (32
original, 17 from \path{L14.lean}, 27 from \path{L15Classical.lean}), the tampering tests are all
rejected, and a scan of both files for \lean{sorry}, \lean{admit}, \lean{native\_decide},
\lean{axiom}, \lean{set\_option}, \lean{macro} and \lean{notation} finds no occurrence outside
comments (the gate's own token check also passes).

\paragraph{Layer 14: statements in Mathlib's language and non-vacuity.} \path{Baranyai/L14.lean}
restates 8 theorems with Mathlib-only types and proves each from the library through bridges:
\decl{baranyai}, \decl{baranyai'}, \decl{complete_graph_one_factorisation}, \decl{baranyai_four_two},
\decl{Baranyai.baranyai_of_card}, \decl{Baranyai.baranyai_finset}, \decl{Baranyai.baranyai_partition}
and the converse \decl{Baranyai.dvd_card_of_factorisation}. The colouring is replaced by an injective
family $P$ of $\binom{n-1}{k-1}$ \mdecl{Finpartition}s of the vertex set into $k$-blocks, each with
$n/k$ blocks, such that every $k$-subset is a block of exactly one $P_i$ (\decl{l14_baranyai});
Theorem~\ref{thm:partition} becomes one \mdecl{Finpartition} of the set of $k$-subsets with
$\binom{n-1}{k-1}$ parts, each a partition of the vertices into $n/k$ blocks
(\decl{l14_baranyai_partition}); the case $k=2$ becomes $n-1$ distinct subgraphs of the complete
\lean{SimpleGraph} satisfying \mdecl{SimpleGraph.Subgraph.IsPerfectMatching}, each edge in exactly
one (\decl{l14_complete_graph_one_factorisation}); the converse says that a \mdecl{Finpartition} of
the vertices into $k$-blocks forces $k\mid|\alpha|$ (\decl{l14_dvd_card_of_factorisation}). A
mechanical check of the elaborated types found no constant of the library in any of them. Each
theorem has a witness theorem that proves the hypotheses for explicit values and proves the
conclusion non-degenerate: for $(n,k)=(6,3)$, ten pairwise distinct partitions with two triples each
(\decl{l14w_baranyai}, and the same for \lean{Fin 2 × Fin 3}); for $(6,2)$, five distinct partitions
with three pairs each; for $K_6$, five distinct spanning perfect matchings; for $K_4$, three distinct
perfect matchings, and a hand-written factorisation of $K_4$ checked by \lean{decide}
(\decl{l14w_baranyai_explicit}); for $\{0,2,4,6\}\subseteq\NN$, three distinct partitions into
pairs; for the partition form at $(6,3)$, ten classes, each with two distinct disjoint triples; and
for the converse, a partition of \lean{Fin 6} into two triples, and the consequence that \lean{Fin 5}
has no partition into pairs. The other declarations of Table~\ref{tab:thms} (the invariant, the
demand, the arithmetic and counting lemmas, the induction step, \decl{Baranyai.factorisation_congr}
and the partition lemmas) are proof internals or are subsumed by the restated theorems, and were not
restated. No Mathlib-only statement turned out false, and none needed an extra hypothesis; the only
additions are \lean{DecidableEq} instances required by \mdecl{Finpartition}. The restatement
confirms the discussion of $n=0$ in Section~\ref{sec:statement}: the one-\mdecl{Finpartition} form
genuinely needs $0<n$ (for $n=0$, $k=1$ it would have $0$ parts but asserts $\binom{0}{0}=1$), while
the family form \decl{l14_baranyai} holds at $n=0$ without it.

\paragraph{Layer 15: comparison with textbook statements.} Another agent collected three statements
from the literature, with citations (Baranyai~\cite{Baranyai1975}, van Lint and
Wilson~\cite[Thm.~38.1]{vLW}, and the Wikipedia article): Baranyai1, the hypergraph form; Baranyai2,
the parallel-class form and its ``equivalently'' version; and Baranyai3, the general almost-regular
form with its parts (a) $k\mid n$ and (b) $k=2$. An AI model, given only these statements (it saw no
code and not this paper; its transcript was audited), wrote them in Lean with Mathlib notions and
four definitions of its own (\lean{IsPerfectMatchingK}, \lean{IsPartitionOfKSubsets},
\lean{pointDegree}, \lean{IsAlmostRegular}): 8 statements, all with \lean{sorry}, with three
ambiguities marked and no extra hypothesis. A comparer then proved each statement, copied verbatim,
from the development, or classified it. \path{Baranyai/L15Classical.lean} (namespace
\lean{Baranyai.Classical}) contains the result. Of the 8 statements, 6 are covered (equivalent to or
implied by the library), 1 is a new proof of a statement outside the paper, and 1 is out of scope;
0 are narrower in the library than in the text, and 0 writer statements are false as written.
\begin{itemize}
\item \emph{Covered (6):} \decl{Baranyai.Classical.cov_Baranyai1_a} (hypergraph form; the writer's
$0<n$ is not needed), \decl{Baranyai.Classical.cov_Baranyai1_b} ($\binom{n-1}{k-1}n=\binom nk k$, from
\decl{Baranyai.choose_mul}), \decl{Baranyai.Classical.cov_Baranyai2_a} and
\decl{Baranyai.Classical.cov_Baranyai2_b} (parallel classes as \mdecl{Finpartition}s; both readings
of the writer's ambiguity at $n=0$ are covered), \decl{Baranyai.Classical.cov_Baranyai3_a}
(the case $k\mid n$ of the almost-regular form), and \decl{Baranyai.Classical.cov_Baranyai3_b}
($K_n$, $n$ even, as $n-1$ perfect matchings of the complete \lean{SimpleGraph} on any finite type,
including $n=0$).
\item \emph{New proof (1):} \decl{Baranyai.Classical.cov_Baranyai3_a_implication}: for $k\mid n$, an
almost-regular partition into classes of size $n/k$ consists of perfect matchings (every point
degree is $1$). The library has no such statement.
\item \emph{Out of scope (1):} the general almost-regular form (van Lint--Wilson Theorem~38.1:
arbitrary sizes $a_1+\dots+a_t=\binom nk$, $k$ not necessarily dividing $n$). This is the
generalisation that Section~\ref{sec:limits} lists as not formalised. It is kept as the unproved
proposition \decl{Baranyai.Classical.classical_Baranyai3_main}, used and assumed nowhere; only its
instance $k\mid n$, $t=\binom{n-1}{k-1}$, $a_i=n/k$ is proved
(\decl{Baranyai.Classical.cov_Baranyai3_main_part}).
\end{itemize}
The bridges written for the comparison (from colourings to their classes, to \mdecl{Finpartition}s
and to Mathlib perfect matchings, and the point-degree computation) are also new proofs. Each proved
\lean{classical\_\dots} theorem is the writer's text with only \lean{sorry} replaced, checked
byte-for-byte against the writer's file. In two places the library is stronger than the textbook text
($n=0$ is covered; no $0<n$ is needed in Baranyai1), and \decl{baranyai'} and
\decl{Baranyai.dvd_card_of_factorisation} have no textbook counterpart among the statements. The
statement collection, the Lean statements and the comparison were all done by AI agents; no human
has checked the textbook statements against the printed sources (the collector notes that its
wording of van Lint--Wilson Theorem~38.1 is from memory).

\section{Related work}\label{sec:related}

\paragraph{Mathematics.} Baranyai's original proof is in~\cite{Baranyai1975}; \cite{Baranyai1979}
treats edge-colourings of complete hypergraphs more generally. The proof formalised here is Brouwer
and Schrijver's~\cite{BrouwerSchrijver}, as presented in van Lint and Wilson~\cite[Ch.~38]{vLW},
where it is noted (p.~537) that all known proofs use a form or consequence of their Theorem~7.2 or~7.4,
the integrality theorems for flows and circulations. Our rounding step uses one such consequence, Hall's theorem~\cite{Hall}, through the token
construction of Section~\ref{sec:proof}. Mendelsohn and Rosa survey 1-factorisations of complete
graphs~\cite{MendelsohnRosa}, and Peltesohn treated the case $k=3$~\cite{Peltesohn}.

\paragraph{Formalisations.} Our searches were made on 2 October 2026. Within Mathlib (source at
revision \texttt{0df444a3}) the only occurrence of ``Baranyai'' is the \path{docs/1000.yaml} entry
quoted in the introduction, with no declaration recorded~\cite{Mathlib1000}; the theorem is not in
\path{docs/100.yaml}, and Mathlib's \lean{Hypergraph} file defines no matchings or factorisations.
The 1000+ theorems project page for \texttt{Q790236} lists no formalisation~\cite{Plus1000}. On GitHub
we ran the following searches. Code search for \texttt{baranyai} restricted to the Lean language
returned no results. Unrestricted code search for \texttt{baranyai} returned results of which we
screened the first 50 file paths only for the extensions \texttt{.lean}, \texttt{.v}, \texttt{.thy},
\texttt{.ml} and \texttt{.miz}, and none had them. Code search for the phrase ``Baranyai's
theorem'' was screened on its first 20 results (repository and file path): apart from the two entries cited above~\cite{Mathlib1000,
Plus1000}, they were prose, data sets, glossaries and lecture notes, and none was a formal proof.
Code search for ``theorem baranyai'', and for ``hypergraph factorization'' restricted to Lean,
returned nothing, and code search for \texttt{baranyai}
restricted to the file extensions \texttt{.v} (Rocq) and \texttt{.thy} (Isabelle) also returned
nothing. Searches of pull requests and issues of the Mathlib repository for ``Baranyai'' returned
none, and repository search for \texttt{baranyai} returned only unrelated repositories named after
people. Code search within the GitHub mirror of
the Isabelle Archive of Formal Proofs~\cite{AFP}, the HOL~Light repository~\cite{HOLLight}, and the
Rocq repositories \texttt{coq-100-theorems}, \texttt{math-comp} and \texttt{graph-theory} returned
nothing, and the list of famous theorems proven in Rocq~\cite{Rocq100} does not mention the theorem.
For Mizar and for the Lean Zulip chat we relied on web search only, including the query
\texttt{"Baranyai" site:leanprover.zulipchat.com}, which returned no results; the other web
searches returned nothing relevant; we did not search the Mizar Mathematical Library
or the Zulip chat directly. These searches are not exhaustive. As far as we could determine, no
earlier formal proof of Baranyai's theorem was public.

\section{Limitations and further work}\label{sec:limits}

\begin{itemize}
\item \emph{Bridges.} The library statements use a direct $\exists!$ formulation. The version-2
files restate them with \lean{Finpartition} and \lean{SimpleGraph} perfect matchings
(Section~\ref{sec:l1415}); a bridge to Mathlib's \lean{Hypergraph} is still missing, and the bridges
live in verification files, not in the library's own API.
\item \emph{Generalisations.} Baranyai's method gives more: almost-regular colourings when
$k\nmid n$, and factorisations of complete multipartite and non-uniform hypergraphs. None is
formalised. The invariant would have to be generalised. The Layer-15 comparison states the general
almost-regular form (van Lint--Wilson Theorem~38.1) as an unproved proposition and proves only its
divisible instance (Section~\ref{sec:l1415}).
\item \emph{Audit coverage.} The statement-level re-run covers all three library files: mutation,
the contradiction probe, ablation and the linter were applied to \path{Partition.lean} and
\path{Transport.lean} as well, with 93 of 93 mutations rejected. Two theorems yielded no mutants
(the mutator produced no conclusion edit for them), and the two hypotheses of \decl{baranyai'}
remain unclassified because their control copy does not compile; the per-mutation diagnostics were
adjudicated individually only for the first run on \path{Basic.lean}.
\item \emph{Cleanup.} The unused hypothesis of \decl{Baranyai.pascal_dem} could be removed; we left
the frozen commit unchanged so that the line references of this paper remain valid.
\end{itemize}

\section*{Acknowledgements and use of AI}
A large-language-model AI assistant was used to write the Lean proofs and the Python check, to run
the prior-art searches and the audits, and to draft this paper. In version~2, AI agents also wrote
the Layer-14 restatements, collected the textbook statements, wrote their Lean versions without access
to the code or the paper, and carried out the comparison (Section~\ref{sec:l1415}). The author directed the work and
reviewed the paper, and checked each formal statement in Table~\ref{tab:thms} against its informal
counterpart by hand. Every proof is checked by the Lean kernel, and the author takes responsibility
for all content.

\begin{thebibliography}{99}\sloppy
\bibitem{AFP} \emph{Archive of Formal Proofs}, \url{https://www.isa-afp.org/}; GitHub mirror
\url{https://github.com/isabelle-prover/mirror-afp-devel}, accessed 2 October 2026.
\bibitem{Baranyai1975} Zs.~Baranyai, On the factorization of the complete uniform hypergraph, in
A.~Hajnal, R.~Rado and V.~T.~S\'os (eds.), \emph{Infinite and Finite Sets} (Colloq.\ Keszthely,
1973), Vol.~I, Colloquia Mathematica Societatis J\'anos Bolyai~10, North-Holland, Amsterdam, 1975,
91--108.
\bibitem{Baranyai1979} Zs.~Baranyai, The edge-coloring of complete hypergraphs~I, \emph{Journal of
Combinatorial Theory, Series~B} 26 (1979), 276--294.
\bibitem{BrouwerSchrijver} A.~E.~Brouwer and A.~Schrijver, Uniform hypergraphs, in A.~Schrijver
(ed.), \emph{Packing and Covering in Combinatorics}, Mathematical Centre Tracts~106, Mathematisch
Centrum, Amsterdam, 1979, 39--73.
\bibitem{HallLean} A.~Gusakov, B.~Mehta and K.~A.~Miller, Formalizing Hall's marriage theorem in
Lean, arXiv:2101.00127, 2021.
\bibitem{Hall} P.~Hall, On representatives of subsets, \emph{Journal of the London Mathematical
Society} 10 (1935), 26--30.
\bibitem{HOLLight} J.~Harrison, HOL~Light repository, \url{https://github.com/jrh13/hol-light},
accessed 2 October 2026.
\bibitem{Lean4} L.~de Moura and S.~Ullrich, The Lean~4 theorem prover and programming language, in
\emph{Automated Deduction -- CADE~28}, Lecture Notes in Computer Science 12699, Springer, 2021,
625--635.
\bibitem{Mathlib} The mathlib Community, The Lean mathematical library, in \emph{Proceedings of the
9th ACM SIGPLAN International Conference on Certified Programs and Proofs (CPP 2020)}, ACM, 2020,
367--381.
\bibitem{Mathlib1000} The mathlib Community, \path{docs/1000.yaml}, entry \texttt{Q790236}, Mathlib
repository, \url{https://github.com/leanprover-community/mathlib4}, revision \texttt{0df444a3},
accessed 2 October 2026.
\bibitem{MendelsohnRosa} E.~Mendelsohn and A.~Rosa, One-factorizations of the complete graph---a
survey, \emph{Journal of Graph Theory} 9 (1985), 43--65.
\bibitem{Peltesohn} R.~Peltesohn, \emph{Das Turnierproblem f\"ur Spiele zu je dreien}, Dissertation,
Berlin; August Pries, Leipzig, 1936.
\bibitem{Plus1000} \emph{1000+ theorems} project, page \path{_thm/Q790236.md},
\url{https://github.com/1000-plus/1000-plus.github.io}, accessed 2 October 2026.
\bibitem{Rocq100} \emph{Famous theorems proven in Rocq}, \url{https://madiot.fr/rocq100/}, accessed
2 October 2026.
\bibitem{vLW} J.~H.~van Lint and R.~M.~Wilson, \emph{A Course in Combinatorics}, 2nd edition,
Cambridge University Press, 2001.
\end{thebibliography}

\end{document}
